\section{Declaración de uso de inteligencia artificial}
Durante el Sprint 3 el equipo usó herramientas de inteligencia artificial generativa, entre ellas Claude Code y Codex, para escribir gran parte del código, generar pruebas automáticas, revisar errores y preparar este informe. Los integrantes revisaron el código y los resultados, ejecutaron las pruebas automáticas y la integración continua, y realizaron la ronda de pruebas manuales descrita en la sección 10. La IA local Ollama forma parte del sistema SIHA como apoyo para analizar materias; no se usó como herramienta de desarrollo.

\section{Conclusión y próximos pasos}
El Sprint 1 entregó la base del sistema: cuentas, permisos, reglas, catálogo y flujo inicial de expedientes. El Sprint 2 incorporó la lectura del plan, el análisis de créditos y contenidos, la explicación de resultados y la consulta de precedentes, además de sus correcciones posteriores. El Sprint 3 completó el flujo con el certificado SIHA único, la revisión previa, los sílabos oficiales y externos, las recomendaciones por materia, los votos de dos docentes, la decisión administrativa cuando discrepan, el registro manual y los resultados consultables.

La regla quedó aclarada con el docente: si los créditos están bajo el 80~\%, la materia queda bloqueada; si los contenidos están bajo el 80~\%, el agente recomienda no homologar, pero los docentes pueden decidir con su criterio. El agente no sustituye la revisión de los docentes.

El hallazgo H-03, en el que las sugerencias dejaban vacía la materia de destino, ya quedó corregido. Los siguientes pasos son decidir si se recupera H-04, la página de reglas; sincronizar el estado del expediente en MS B cuando MS C lo resuelve (H-05); y mostrar con tildes los nombres del catálogo (H-06). También quedan por evaluar los avisos por correo y cargar los sílabos de Robótica para los niveles 6 a 9.
